\documentclass[12pt]{article}
\usepackage[a3paper, margin=1in]{geometry}
\usepackage{amsmath, amssymb, amsthm, ulem}

\begin{document}

\section*{Domácí úkol 6.2}

Vypracoval: Daniel \textit{"Randál"} Ransdorf \hfill Podpis: \rule{4cm}{0.4pt}

\begin{enumerate}
  \setcounter{enumi}{1}
  \item Když je obraz B roven jádru A. Znamená to, že sloupcový prostor B je roven
    řešení homogenní soustavy $A\vec{x} = \vec{0}$. Nejprve vezměme matici $B^T$,
    která má řádky rovny sloupcům $B$ a 
    proveďme Gaussovu eliminaci, čímž řádkový prostor nezměníme, jen zjednodušíme:

    \[
    \left(\begin{array}{ccccc}
      0&2&0&2&2 \\ 0&1&2&2&0 \\ 0&1&1&0&2
    \end{array}\right)
    \overset{\text{2.ř }+\text{1.ř}}{\sim}
    \left(\begin{array}{ccccc}
      0&2&0&2&2 \\ 0&0&2&1&2 \\ 0&1&1&0&2
    \end{array}\right)
    \overset{\text{3.ř }+\text{1.ř}}{\sim}
    \left(\begin{array}{ccccc}
      0&2&0&2&2 \\ 0&0&2&1&2 \\ 0&0&1&2&1
    \end{array}\right)
    \overset{\text{3.ř }+\text{2.ř}}{\sim}
    \left(\begin{array}{ccccc}
      0&2&0&2&2 \\ 0&0&2&1&2 \\ 0&0&0&0&0
    \end{array}\right)
    \]

    Plyne, že obraz $B$ je roven řádkovému prostoru této zeliminované matice, tedy:

    \[
      Im(B) = Span\left(\left\{
        \left(
          \begin{array}{c}
            0\\2\\0\\2\\2
          \end{array}
        \right),
        \left(
          \begin{array}{c}
            0\\0\\2\\1\\2
          \end{array}
        \right)
      \right\}\right)
    \]

    Podle původní úvahy plyne, že řešení homogenní soustavy s maticí $A$ je:

    \[
      \left\{
        t_1\left(\begin{array}{c}
          0\\2\\0\\2\\2
        \end{array}\right) +
        t_2 \left(\begin{array}{c}
          0\\0\\2\\1\\2
        \end{array}\right)
        \quad
        \mid
        t_1,t_2 \in \mathbb{Z}_3
      \right\}
    \]

    Řešením jsou pěti-složkové aritmetické vektory, tedy matice $A$ musí mít 5 sloupců.
    Řešením je dvojrozměrný prostor, tedy matice $A$ má 3 bázové sloupce a 2 volné.
    Hodnost $A$ tedy musí být 3, a musí tedy mít alespoň tři řádky.

    Víme jak vypadá řešení homogenní soustavy.
    Zapišme tedy Gauss-Jordan zeliminovaný tvar homogenní soustavy k $A$,
    ze které můžeme zpětně homogenní soustavu získat "vynulováním" pravé strany:

    \[
      \left(\begin{array}{ccccc|cc}
        1 & 0 & 0 & 0 & 0 & 0 & +0 \\
        0 & 1 & 0 & 0 & 0 & 2t_1 & +0 \\
        0 & 0 & 1 & 0 & 0 & 0 & +2t_2 \\ \hline
        0 & 0 & 0 & 1 & 0 & 2t_1 & +t_2 \\
        0 & 0 & 0 & 0 & 1 & 2t_1 & +2t_2 \\
      \end{array}\right)
      \overset{\text{2.ř }+\text{4.ř}}{\;\sim\;}
      \left(\begin{array}{ccccc|c}
        1 & 0 & 0 & 0 & 0 & 0  \\
        0 & 1 & 0 & 1 & 0 & t_1+t_2 \\
        0 & 0 & 1 & 0 & 0 & 2t_2 \\ \hline
        0 & 0 & 0 & 1 & 0 & 2t_1+t_2 \\
        0 & 0 & 0 & 0 & 1 & 2t_1+2t_2 \\
      \end{array}\right)
      \overset{\text{2.ř }+\text{5.ř}}{\;\sim\;}
      \left(\begin{array}{ccccc|c}
        1 & 0 & 0 & 0 & 0 & 0  \\
        0 & 1 & 0 & 1 & 1 & 0 \\
        0 & 0 & 1 & 0 & 0 & 2t_2 \\ \hline
        0 & 0 & 0 & 1 & 0 & 2t_1+t_2 \\
        0 & 0 & 0 & 0 & 1 & 2t_1+2t_2 \\
      \end{array}\right)
    \]
    \[
      \overset{\text{3.ř }+\text{5.ř}}{\;\sim\;}
      \left(\begin{array}{ccccc|c}
        1 & 0 & 0 & 0 & 0 & 0  \\
        0 & 1 & 0 & 1 & 1 & 0 \\
        0 & 0 & 1 & 0 & 1 & 2t_1 + t_2 \\ \hline
        0 & 0 & 0 & 1 & 0 & 2t_1+t_2 \\
        0 & 0 & 0 & 0 & 1 & 2t_1+2t_2 \\
      \end{array}\right)
      \overset{\text{3.ř }+2\cdot\text{4.ř}}{\;\sim\;}
      \left(\begin{array}{ccccc|c}
        1 & 0 & 0 & 0 & 0 & 0  \\
        0 & 1 & 0 & 1 & 1 & 0 \\
        0 & 0 & 1 & 2 & 1 & 0 \\ \hline
        0 & 0 & 0 & 1 & 0 & 2t_1+t_2 \\
        0 & 0 & 0 & 0 & 1 & 2t_1+2t_2 \\
      \end{array}\right)
    \]

    Zapišme nalezenou matici $A$:

    \[
      A = \underline{\underline{\left(\begin{array}{ccccc}
        1 & 0 & 0 & 0 & 0 \\
        0 & 1 & 0 & 1 & 1 \\
        0 & 0 & 1 & 2 & 1 \\
      \end{array}\right)}}
    \]

    Ukažme pro jistotu, že matice jádro $A$ skutečně obsahuje všechny prvky obrazu $B$.
    Tedy že zobrazení $f_A\circ f_B$ zobrazuje všechno do nulového vektoru.
    Protože zobrazení $f_B$ všechno zobrazuje do svého obrazu, ale onoho obraz je jádro $A$,
    tedy $f_A\circ f_B$ musí vše zobrazit do nulového vektoru.
    Tedy součin $A\cdot B$ musí být nulová matice, jinak bychom triviálně nalezli
    vektor $\vec{x}$, který nesplňuje $A\cdot B \cdot \vec{x} = \vec{0}$.

    \[
      AB = \left(\begin{array}{ccccc}
        1 & 0 & 0 & 0 & 0 \\
        0 & 1 & 0 & 1 & 1 \\
        0 & 0 & 1 & 2 & 1 \\
      \end{array}\right) \cdot \left(\begin{array}{ccccc}
        0 & 0 & 0 \\ 2 & 1 & 1 \\ 0 & 2 & 1 \\ 2 & 2 & 0 \\ 2 & 0 & 2
      \end{array}\right) = \left(\begin{array}{ccc}
        0 & 0 & 0 \\ 0 & 0 & 0 \\ 0 & 0 & 0
      \end{array}\right)
    \]

    Kontrola splněna.

\end{enumerate}

\end{document}
